\section{How We Measured}
\label{sec:purse}

To know whether agents spend well, we need tasks where the best split of the money is known. Our benchmark, \bench{},
builds such tasks. Each project is a set of jobs that share one budget. Every job has an exact answer, and every job has
been tried on every option in advance, so the best split for any budget can be computed. The score then depends on how
the money is split, and the gap to the best split shows how much a strategy left on the table.

\subsection{Jobs and projects}
A project has eight jobs, two from each of four job types (Table~\ref{tab:families}), and one budget $B$. Each job has
a visible level from 1 to 5 (its stated size), a hidden difficulty (low, mid or high), a value of 1, 2 or 4 points
chosen independently of difficulty, a standard prompt and an exact checker. A strategy reads every job with its level
and value and answers with one JSON object. The score is the share of points answered correctly. Jobs are generated
from random seeds, so they cannot be looked up, and \sys{}'s web tools are turned off in every run.

\begin{table}[!htb]
\centering
\caption{\bench{}'s four job types.}
\label{tab:families}
\footnotesize
\begin{tabular}{@{}p{2.2cm}p{4.0cm}p{5.2cm}p{3.0cm}@{}}
\toprule
\tablehead Job type & Set by the level & Set by the hidden difficulty & Can be bought from \\
\midrule
Arithmetic & operators (2 to 22), digits (2 to 8) & multiplications, digit sizes (carries) & \texttt{calc} (levels 1 to 4) \\
Logic grids & houses, attributes & guesses a constraint solver must make (0, 1, 2+) & \texttt{csp}; a free checker \\
Aggregation & records (40 to 2{,}500) & fields in the condition (1, 2, or 2 plus near misses) & no one \\
String pipelines & operations (2 to 10), length & alphabet, operations that depend on position & no one \\
\bottomrule
\end{tabular}
\end{table}

For the second round (\S\ref{sec:evaluation}), the hidden difficulties were spread further apart, so that knowing them
would be worth more. In arithmetic, the hidden difficulty also changes the numbers' digits (one fewer, the same, or two
more) and the share of multiplications (10\%, 35\% or 70\%). In aggregation, it sets how many records match (a handful,
about 1 in 20, or about 1 in 7) and triples the near misses. In strings, it halves, keeps or doubles the starting
string. At high difficulty all but one operation depends on position, and at low difficulty the Caesar shift is
dropped.

\subsection{How much of the difficulty is visible}
With probability $\rho$, a job's hidden difficulty is the one its visible level suggests (levels 1 and 2 low, 3 mid, 4
and 5 high). Otherwise it is random. At $\rho = 1$ the level tells the whole story, so a plan made up front with the
right statistics loses nothing to a strategy that adapts. At $\rho = 0$, difficulty can only be learned by trying. The
value of knowing the hidden difficulty,
\[
\mathrm{VOI} = \mathrm{score}(\text{oracle knowing the difficulty}) - \mathrm{score}(\text{oracle knowing only the levels}),
\]
is the most that any adaptive strategy can gain over the best up-front plan, and we report it for every setting. Each
strategy gets the statistics it would measure on projects with its own $\rho$, which are each difficulty's success
rates weighted by $P(\text{difficulty} \mid \text{level}) = \rho\,[\text{difficulty matches the level}] + (1 - \rho)/3$.

\subsection{The market}
\label{sec:purse-market}
The options are models at different prices (a cheap and an expensive reasoning model in the first round, plus a third,
\texttt{gpt-5-nano}, in the second) and outside agents reached over A2A. Two outside agents are exact programs.
\texttt{calc} does exact arithmetic (it is \sys{}'s example agent, and reads expressions of up to 200 characters, so it
refuses level 5), and \texttt{csp} solves logic grids. The second round adds \texttt{gridfast}, which sells the same
grids at a quarter of \texttt{csp}'s price and gets 40\% of them wrong (it swaps one attribute between two houses), so
that the cheapest seller is not always a good one. Outside agents charge a price per message, tested at
$c_\mathrm{ref} \times \{0.25, 1, 4, 16\}$, where $c_\mathrm{ref}$ is the cheap model's median cost on the jobs they can
do. A free \texttt{grid\_check} tool says whether a grid solution meets every clue, up to three times per puzzle
(checking is cheaper than solving, as with unit tests).

\subsection{Recorded results, replay and the oracle}
Before any strategy runs, every job is sent with its standard prompt to every option that can do it (once in the first
round, three times in the second), through a \sys{} hub (one hire, one message, one fire). We call this
\emph{calibration}. The \emph{results table} records each answer, whether it is right, and what it cost. A replayed
call returns one of the job's recorded attempts, and each seed starts from a different one. The table has three uses.
First, it gives the informed strategies their statistics, taken from the other half of the jobs, so a project never
sees statistics that include its own jobs. Second, it defines the \emph{oracle}. Oracle-full picks, for each job, the
option that gets the most points under the budget, using each job's recorded success and cost (a multiple-choice
knapsack, as in $R^3$-Bench~\cite{wang2026r3bench}). No strategy can beat it. Third, it lets any strategy that sends
standard prompts be scored by \emph{replay} on recorded results instead of new model calls, as RouterBench does for
routers~\cite{hu2024routerbench}. Each project's budget is $B = \tau\,C_\mathrm{full}$, where $C_\mathrm{full}$ is what
oracle-full spends at its best score when money is unlimited. So $\tau = 1$ is just enough money for the best answers,
and $\tau = 0.5$ is half of that.

\subsection{Strategies}
\label{sec:purse-allocators}
The main comparison crosses who splits the money (a language model or a solver) with when (up front, or during the
task). All four get the same project text, prices, statistics and budget (Table~\ref{tab:twobytwo}).

\begin{table}[!htb]
\centering
\caption{The four strategies in the main comparison.}
\label{tab:twobytwo}
\footnotesize
\begin{tabular}{@{}p{1.6cm}p{6.4cm}p{6.4cm}@{}}
\toprule
\tablehead  & Up front & During the task \\
\midrule
LLM & \textbf{LLM split.} The CEO plans one hire per job in plan mode, and the plan runs unchanged (like ZEBRA's LLM
baseline). & \textbf{\sys{}.} The CEO hires, asks, fires, rehires and moves up to better models as results come in. \\
Solver & \textbf{Knapsack.} Picks options to get the most expected points under the budget, using the statistics. &
\textbf{Cascade.} Takes the cheapest option expected to succeed, and moves up a tier when a check fails or two answers
disagree. \\
\bottomrule
\end{tabular}
\end{table}

Around them are oracle-full, fixed strategies (every job on one model; a hybrid that picks by visible level, in the
style of AgentCARD~\cite{jiang2026specialize}; an equal split; a random choice), and \sys{} variants (a CEO that names
models, as in v1, and a CEO that sees prices but no statistics).

\subsection{Validity checks}
\label{sec:checks}
Before any strategy runs, seven checks on the results table make sure \bench{} really measures how the money is split
(Table~\ref{tab:checks}). We tested them on made-up tables built to pass, and on one table per check built to fail
exactly that check. Check V8 (replay agrees with live runs) needs live runs.

Both tables behind \S\ref{sec:evaluation} fail three of the seven, and not the same three. On the first table, V1
failed narrowly (the oracle beat the best fixed strategy by 9.7 points, not 10). V4 failed because the expensive model
had room to improve in only one job type (long aggregation jobs), not two. V5 failed outright. Knowing each job's
hidden difficulty was worth 0.009 of the points at $\rho = 0$, against the 0.08 the check asks for. We built the second
table to fix these, with a third model, a cheap but unreliable grid seller, more spread-out difficulties and three
attempts per job. V1 now holds (10.7 points). V5 still fails, and by more (0.002). The wider spread did separate the
cheap models (\texttt{gpt-6-luna} answers 100\% of low-difficulty aggregation jobs and 43\% of mid ones). But knowing
the hidden difficulty is worth between $-0.007$ and 0.002 of the points at every budget we tried ($\tau$ from 0.25 to
2), and 0.01 to 0.02 with the expensive model removed. With only four jobs per group, each difficulty's rates rest on
twelve attempts, which are as noisy as the differences they are meant to show. V3 now fails because
\texttt{gpt-5-nano} is not simply a weaker model. It beats \texttt{gpt-6-luna} on arithmetic (53\% against 25\% at
level 5) and loses to it on aggregation and strings, so the three models cross over by job type instead of lining up,
and on logic grids the expensive model never leads the cheap one by 20 points. V4 fails as before. So neither table can
show what adapting during the task is worth, and the during-the-task column of Table~\ref{tab:twobytwo} cannot be
tested against the up-front one. Normally a failed check means a job generator is retuned and its jobs are measured
again. We report both tables as they are, and a table that passes V5 is on the roadmap (\S\ref{sec:conclusion}).

\begin{table}[!htb]
\centering
\caption{Validity checks on the results table, run before any strategy, and how the two tables behind
\S\ref{sec:evaluation} do.}
\label{tab:checks}
\footnotesize
\begin{tabular}{@{}p{3.0cm}p{6.7cm}p{2.4cm}p{2.4cm}@{}}
\toprule
\tablehead Check & Holds when & First table & Second table \\
\midrule
V1 Splitting matters & the oracle beats the best fixed strategy by $\geq$ 10 points at $\tau = 0.5$ & \textbf{fails} (9.7) & holds (10.7) \\
V2 The budget is tight & the expensive model alone covers $<$ 60\% of the points at $\tau = 0.5$ & holds (59.6\%) & holds (51.7\%) \\
V3 Models separate & for each job type, a level where the cheap model succeeds $\geq$ 80\% (and a small model, if present, $\leq$ 40\%), and one where the expensive model leads the cheap one by $\geq$ 20 points & holds (4 of 4 types) & \textbf{fails} (0 of 4) \\
V4 Room to improve & the expensive model succeeds $\leq$ 60\% at level 5 in $\geq$ 2 job types & \textbf{fails} (1 type) & \textbf{fails} (1 type) \\
V5 Hidden difficulty matters & VOI $\leq$ 0.02 at $\rho = 1$ and $\geq$ 0.08 at $\rho = 0$ & \textbf{fails} (0.000; 0.009) & \textbf{fails} (0.000; 0.002) \\
V6 Buying responds to price & the oracle's buying falls $\geq$ 40 points across the price sweep & holds (53 points) & holds (56 points) \\
V7 No shortcut & the most common answer per job type scores $\leq$ 10\% & holds (4.3\%) & holds (5.6\%) \\
\bottomrule
\end{tabular}
\end{table}

\FloatBarrier  % the validity checks table stays in this section
